\chapter{References, links and bibliographies}
\label{ch:references}

\section{Cross-references}\label{cross-references}

Give anything numbered a label, then refer to it. The number updates itself if things move.

\begin{lstlisting}[language=latex]
\section{Method}\label{sec:method}
See Section~\ref{sec:method} and Equation~\eqref{eq:energy}.
\end{lstlisting}

The \texttt{\textasciitilde{}} is a non-breaking space, which keeps "Section" and its number on one line. A prefix such as \texttt{sec:}, \texttt{eq:}, \texttt{tab:} or \texttt{fig:} is a helpful habit. The \texttt{cleveref} package can add the word for you: \texttt{\textbackslash{}cref\{sec:method\}}.

\section{Citations}\label{citations}

Keep references in a separate \texttt{.bib} file, and cite them by key:

\textbf{refs.bib}

\lstinputlisting{../static/examples/report/refs.bib}

\begin{lstlisting}[language=latex]
\usepackage[round]{natbib}
...
As \citet{clarke2001} shows, this holds \citep{clarke2001}.
...
\bibliographystyle{plainnat}
\bibliography{refs}
\end{lstlisting}

The bibliography tool picks up only the items you cite. Zotero can export a \texttt{.bib} file, and its Better BibTeX add-on keeps that file updated. \texttt{biblatex} with the \texttt{biber} tool is a more flexible alternative, while journals often supply their own \texttt{.bst} style for BibTeX.

\section{A complete small report}\label{a-complete-small-report}

This pulls together everything so far: sections, equations, a table, a chart, cross-references, a bibliography and links.

\textbf{main.tex}

\lstinputlisting{../static/examples/report/main.tex}

It compiles to these two pages. The links are coloured teal, and every number comes from a label.

\begin{figure}
\centering
\includegraphics[width=\linewidth]{img/report-page-1.png}
\caption{Page 1}
\end{figure}

\begin{figure}
\centering
\includegraphics[width=\linewidth]{img/report-page-2.png}
\caption{Page 2}
\end{figure}

